\begin{Exercise}[label={ex:vacanza}]

(TdE Gennaio 2019) Non vado in vacanza se almeno due dei tre seguenti predicati non sono soddisfatti: (A) ho un gruzzoletto da parte, (B) ho almeno due settimane in cui posso non studiare, (C) ho finito gli esami a luglio.

		\Question Scrivere la tabella di verit\`a per il predicato ``Non vado in vacanza''  attraverso i predicati semplici A, B e C.
		
		\Question La traduzione di tale condizione con $((!A \; || \; !B) \; \&\& \; (!A \; || \; !C) \; \&\& \; (!B \; || \; !C))$ \`e corretta?
\end{Exercise}

\begin{Answer}[ref={ex:vacanza}]

\begin{enumerate}
	\item Due dei tre predicati devono essere falsi perch\'e io non vada in vacanza, quindi la traduzione della precedente condizione \'e:
	
			(! A \&\& ! B) || (! A \&\& ! C) || (! B \&\& ! C)

			in tabella della verit\`a:
			
			\begin{table}[H]
\centering
\begin{tabular}{c|c|c|c}
A & B & C & Non vado in vacanza \\ \hline
0 & 0 & 0 & 1                   \\
0 & 0 & 1 & 1                   \\
0 & 1 & 0 & 1                   \\
0 & 1 & 1 & 0                   \\
1 & 0 & 0 & 1                   \\
1 & 0 & 1 & 0                   \\
1 & 1 & 0 & 0                   \\
1 & 1 & 1 & 0                  
\end{tabular}
\end{table}

	\item Per la verifica che il predicato \'e effettivamente corretto possiamo utilizzare le leggi di de Morgan per verificare l'equivalenza con la soluzione del punto 1. In alternativa si pu\`o scrivere la tabella di verit\`a e verificare che rifletta la condizione richiesta.

\end{enumerate}

\end{Answer}